% TEMPLATE-MILC-v3.tex - plantilla maestra UNICA de las guias MILC v3.
%
% La usan las tres familias de guias, cada una con su driver:
%   · Grados 6-11  -> scripts/build-guias-g11.py
%   · Bebras       -> scripts/build-guias-bebras.py
%   · Semillero    -> scripts/build-guias-semillero.py
%
% Antes habia tres copias casi identicas (~400 lineas iguales, 6 distintas):
% cada mejora habia que aplicarla tres veces, y en la practica no se hacia
% (el motor de recursos vivia solo en dos de ellas). Lo que varia entre
% familias esta parametrizado en 6 placeholders que cada driver llena
% (se nombran aqui SIN su sintaxis real: el builder reemplaza por texto plano,
% tambien dentro de los comentarios, y un valor multilinea rompe el preambulo):
%
%   HEADER_IZQ        encabezado izquierdo de pagina
%   HEADER_DER        encabezado derecho de pagina
%   PORTADA_ETIQUETA  rotulo sobre el numero grande (GRADO/MOMENTO/MODULO)
%   PORTADA_SERIE     linea de serie de la portada
%   PORTADA_ID        identificador de la guia en la portada
%   PORTADA_CIRCULO   el circulito de grado (°), o vacio si no aplica
%
% El resto de claves las documenta content/guias/_SCHEMA.md. El script valida
% que no queden placeholders sin reemplazar antes de escribir el archivo final.

\documentclass[11pt,letterpaper]{article}
\usepackage[margin=1.75cm]{geometry}
\usepackage[table]{xcolor}
\usepackage{fontspec}
\usepackage[spanish]{babel}
\usepackage{tikz}
% Librerías TikZ para los diagramas de las guías (bloque `recursos:`):
%  · babel  → babel en español vuelve activos caracteres como «>», y TikZ los usa
%             en sus opciones (>=stealth, ->). Sin esta librería, cualquier
%             diagrama con flechas revienta con «\language@active@arg> has an
%             extra }». Debe ir ANTES de que se dibuje nada.
%  · positioning        → sintaxis «below left=2pt and 3pt».
%  · decorations.pathreplacing → llaves ({brace}) para señalar rangos.
\usetikzlibrary{babel,positioning,decorations.pathreplacing}
\usepackage{graphicx}
% Circuitos: el trabajo de electrónica se hace hoy con micro:bit y sus sensores
% integrados (MakeCode, sin cableado), así que no se necesitan esquemáticos.
% Si algún día entran montajes con protoboard, basta descomentar esta línea:
% los diagramas .tex/.tikz declarados en `recursos:` ya se incrustan solos.
% \usepackage{circuitikz}
\usepackage[most]{tcolorbox}
\usepackage{tabularx}
\usepackage{array}
\usepackage{fancyhdr}
\usepackage{titlesec}
\usepackage{ragged2e}
\usepackage{enumitem}
\usepackage{microtype}

% Helvetica Neue no trae la flecha U+2192, asi que cada «→» escrito literalmente
% en el YAML se compilaba a NADA, en silencio (462 flechas en 61 guias: rutas del
% tipo «paleta → tipografia → lockup» salian rotas). Mapeamos el caracter a su
% version matematica, que si tiene glifo. Asi el autor puede seguir escribiendo
% «→» comodamente en el YAML.
\usepackage{newunicodechar}
\newunicodechar{→}{\ensuremath{\rightarrow}}

\setmainfont{Helvetica Neue}
\setsansfont{Helvetica Neue}
\setlength{\parindent}{0pt}
\setlength{\parskip}{6pt}
\hyphenpenalty=200
\exhyphenpenalty=200
\pretolerance=400
\tolerance=2000
\setlength{\emergencystretch}{1em}
\renewcommand{\arraystretch}{1.25}
\setlist{leftmargin=5mm,itemsep=1mm,topsep=1mm}
\newcolumntype{Y}{>{\RaggedRight\arraybackslash}X}

\definecolor{milcMagenta}{HTML}{E6007E}
\definecolor{milcVino}{HTML}{5A0038}
\definecolor{milcTurquesa}{HTML}{0093A5}
\definecolor{milcVerde}{HTML}{58A923}
\definecolor{milcMostaza}{HTML}{E5B400}
\definecolor{milcOcre}{HTML}{B86600}
\definecolor{milcNegro}{HTML}{241816}
\definecolor{milcGris}{HTML}{F7F7F4}
\definecolor{milcLinea}{HTML}{E8E2DD}
\definecolor{milcGradePrimary}{HTML}{0066FF}
\definecolor{milcGradeDark}{HTML}{003D99}
\definecolor{milcGradeSoft}{HTML}{D6E8FF}
\definecolor{milcGradeAccent}{HTML}{CFE7FF}
\definecolor{milcGradeText}{HTML}{FFFFFF}
\color{milcNegro}

\pagestyle{fancy}
\fancyhf{}
\lhead{\small Tecnología e Informática · Grado 6}
\rhead{\small Guía 4 · MILC v3}
\cfoot{\small\thepage}
\renewcommand{\headrulewidth}{0pt}

\newtcolorbox{titlebox}[1]{
  enhanced,colback=#1,colframe=#1,arc=7mm,boxrule=0pt,
  left=7mm,right=7mm,top=3mm,bottom=3mm,
  before skip=8pt,after skip=8pt,
  fontupper=\sffamily\bfseries\Large\color{white}
}

\newtcolorbox{softbox}[3]{
  enhanced,colback=#2,colframe=#1,borderline west={4pt}{0pt}{#1},
  arc=2mm,boxrule=.4pt,left=4mm,right=4mm,top=3mm,bottom=3mm,
  before skip=6pt,after skip=8pt,title={#3},coltitle=white,
  colbacktitle=#1,fonttitle=\sffamily\bfseries,fontupper=\RaggedRight,
  attach boxed title to top left={xshift=3mm,yshift=-2mm},
  boxed title style={arc=2mm, boxrule=0pt}
}

\newtcolorbox{drawbox}[2]{
  enhanced,colback=white,colframe=#1,arc=4mm,boxrule=1pt,
  left=5mm,right=5mm,top=4mm,bottom=4mm,before skip=8pt,after skip=8pt,
  title={#2},coltitle=white,colbacktitle=#1,fonttitle=\sffamily\bfseries,
  fontupper=\RaggedRight,
  attach boxed title to top left={xshift=4mm,yshift=-2mm},
  boxed title style={arc=3mm, boxrule=0pt}
}

\newtcolorbox{infoband}[1]{
  enhanced,colback=#1!8,colframe=#1!50,arc=2mm,boxrule=.5pt,
  left=4mm,right=4mm,top=2mm,bottom=2mm,before skip=4pt,after skip=4pt,
  fontupper=\small\RaggedRight
}

% Puente narrativo entre secciones (contrato editorial Conéctate).
% Da continuidad: "Acabas de X. Ahora vas a Y porque Z."
% Visualmente: banda izquierda en color del grado, texto en itálica pequeña.
\newtcolorbox{puentebox}{
  enhanced,colback=milcGradeSoft,colframe=milcGradeDark,
  borderline west={3pt}{0pt}{milcGradeDark},
  arc=0mm,boxrule=0pt,left=5mm,right=4mm,top=2.5mm,bottom=2.5mm,
  before skip=4pt,after skip=8pt,
  fontupper=\itshape\small\color{milcNegro}\RaggedRight
}
% La flecha va en modo matematico a proposito: Helvetica Neue NO tiene el glifo
% U+2192, asi que \textrightarrow se compilaba a NADA («Missing character: There
% is no → in font Helvetica Neue») y el puente salia sin flecha. En modo
% matematico la flecha viene de la fuente matematica y si se dibuja.
\newcommand{\puente}[1]{%
  \begin{puentebox}%
    \textcolor{milcGradeDark}{$\rightarrow$}\hspace{2mm}#1%
  \end{puentebox}%
}

\newcommand{\linead}[1]{\noindent\color{milcLinea}\rule{#1}{0.5pt}\color{milcNegro}}
\newcommand{\respuesta}[1]{\par\vspace{2mm}\foreach \n in {1,...,#1}{\noindent\color{milcLinea}\rule{\linewidth}{0.5pt}\par\vspace{5mm}}\color{milcNegro}}
\newcommand{\checkbox}{\textcolor{milcGradeDark}{\fbox{\phantom{x}}}\hspace{5pt}}

\newcommand{\phasecard}[4]{
\begin{tcolorbox}[enhanced,colback=#3,colframe=#2,arc=3mm,boxrule=.5pt,height=3.05cm,valign=center,halign=center,left=1.5mm,right=1.5mm,top=2mm,bottom=2mm]
{\sffamily\bfseries\fontsize{22}{24}\selectfont\textcolor{#2}{#1}}\par
{\sffamily\bfseries\small #4}
\end{tcolorbox}
}

% ── Recursos visuales de la guía ──────────────────────────────────────
% Declarados en el bloque `recursos:` del YAML y emitidos por el builder.
% \guiaFigura   → imagen rasterizada o PDF (foto, carta celeste, montaje).
% \guiaDiagrama → fuente TikZ/circuitikz (.tex) incrustada como vector:
%                 es la vía para circuitos y esquemáticos editables.
% #1 = ruta relativa al .tex · #2 = pie (puede ir vacío)
\newcommand{\guiaFigura}[2]{%
\begin{center}
\includegraphics[width=0.88\linewidth,height=0.38\textheight,keepaspectratio]{#1}\\[1.5mm]
{\footnotesize\itshape\textcolor{milcNegro!75}{#2}}
\end{center}
}

\newcommand{\guiaDiagrama}[2]{%
\begin{center}
\input{#1}\\[1.5mm]
{\footnotesize\itshape\textcolor{milcNegro!75}{#2}}
\end{center}
}

\begin{document}
\thispagestyle{empty}
\begin{tikzpicture}[remember picture,overlay]
  \fill[white] (current page.south west) rectangle (current page.north east);
  \path[fill=milcGradePrimary,rounded corners=16mm]
    ([xshift=.55cm,yshift=-.55cm]current page.north west) rectangle
    ([xshift=-.55cm,yshift=3.65cm]current page.south east);

  \node[anchor=north west,text=milcGradeText] at ([xshift=2.0cm,yshift=-1.85cm]current page.north west)
    {\sffamily\bfseries\fontsize{14}{16}\selectfont GRADO};
  \node[anchor=north west,text=milcGradeText,opacity=.82] at ([xshift=2.0cm,yshift=-2.75cm]current page.north west)
    {\sffamily\bfseries\fontsize{9}{11}\selectfont SERIE GUÍAS · TECNOLOGÍA E INFORMÁTICA};

  \begin{scope}[shift={([xshift=-3.05cm,yshift=-2.55cm]current page.north east)}]
    \fill[white,opacity=.80] (-.62,-.52) rectangle (.62,.52);
    \draw[milcGradeText,line width=1pt] (-.62,-.52) rectangle (.62,.52);
    \fill[milcGradeDark] (-.42,-.38) rectangle (-.22,.06);
    \fill[milcGradeAccent] (-.10,-.38) rectangle (.10,.24);
    \fill[milcGradePrimary!60!black] (.22,-.38) rectangle (.42,.38);
  \end{scope}

  \node[anchor=west,text=milcGradeText] at ([xshift=1.9cm,yshift=-12.85cm]current page.north west)
    {\sffamily\bfseries\fontsize{124}{124}\selectfont 6};
  % El circulito de grado (°) solo aplica a las guias de grado («11°»); en
  % Bebras y Semillero el numero grande es un momento/modulo, no un grado.
  \draw[milcGradeText,line width=1.7pt]
    ([xshift=4.52cm,yshift=-11.68cm]current page.north west) circle (.13cm);

  \node[anchor=west,text=milcGradeText,text width=8.7cm,align=left] at ([xshift=10.35cm,yshift=-11.6cm]current page.north west)
    {
     {\sffamily\bfseries\fontsize{19}{22}\selectfont Netiqueta ---\\los 10 acuerdos\\para llevarte bien\\en internet}\\[4mm]
     {\sffamily\bfseries\fontsize{23}{26}\selectfont Guía 4-6-TIC}\\[2mm]
     {\sffamily\fontsize{13}{16}\selectfont Período 1 · Comunicación e identidad digital}\\[5mm]
     {\sffamily\bfseries\fontsize{11}{13}\selectfont Institución Educativa Sor María Juliana}\\[1mm]
     {\sffamily\fontsize{10}{12}\selectfont Autor: Dr. Álvaro Cárdenas Orozco}
    };

  \path[fill=milcGradeSoft,rounded corners=7mm]
    ([xshift=1.35cm,yshift=.85cm]current page.south west) rectangle
    ([xshift=-1.35cm,yshift=4.25cm]current page.south east);
  \draw[milcGradeDark,line width=.65pt,rounded corners=7mm]
    ([xshift=1.35cm,yshift=.85cm]current page.south west) rectangle
    ([xshift=-1.35cm,yshift=4.25cm]current page.south east);
  \node[anchor=center,text=milcNegro,text width=17.0cm,align=center] at ([yshift=2.55cm]current page.south)
    {\sffamily\fontsize{10}{12}\selectfont Metodología MILC · Apertura ancestral · Pensamiento computacional · Triángulo Dussel-Estoicismo-Floridi\\
    \textcolor{milcGradeDark}{\bfseries Producto:} Un \textbf{cartel personal de los 10 acuerdos de netiqueta} dibujado en hoja completa del cuaderno. Cada acuerdo va con \textbf{un ejemplo de cumplirlo} y \textbf{un ejemplo de romperlo} con tus propias palabras. Más \textbf{una promesa firmada} al final: \emph{``Me comprometo a respetar estos 10 acuerdos en todos los chats, juegos y redes en que esté.''} Lo muestras antes de salir.};
\end{tikzpicture}
\mbox{}
\newpage

\section*{Apertura: Netiqueta --- los 10 acuerdos para llevarte bien en internet}

\begin{softbox}{milcGradeDark}{milcGradeSoft}{Saber ancestral}
\textbf{Hay reglas que no están escritas pero todos las saben.} En tu casa hay normas que nadie te enseñó por escrito: no se grita en la mesa, no se interrumpe al que está hablando, se saluda al entrar. Esas reglas tu mamá o tu abuela no te las leyeron de un libro: las aprendiste viendo. En las cuadras del centro de Cartago pasaba parecido: cuando llegaba el pregonero, los niños se callaban. Era una regla no escrita pero todos la respetaban porque así el mensaje llegaba completo. Si alguien gritaba o se reía, otra persona le decía: \emph{``Cállese mijo, déjelo hablar''}. \textbf{Esas reglas no escritas se llaman acuerdos de convivencia}. En internet también existen --- pero como nadie te las enseña por verlas, te las tienen que enseñar a propósito. Se llaman \textbf{netiqueta}: la etiqueta (las reglas de cortesía) en la red.
\end{softbox}

\begin{softbox}{milcTurquesa}{milcGris}{Saber contemporáneo}
\textbf{Netiqueta} viene de juntar dos palabras: \emph{net} (red en inglés) + \emph{etiqueta} (reglas de cortesía). Son los acuerdos para tratarte bien con otros cuando estás en chats, juegos, comentarios, correos y grupos. \textbf{No son una ley:} nadie te lleva a la cárcel si las rompes. Son algo más fuerte: \textbf{son la forma en que la gente decide si quiere o no juntarse contigo en línea}. Quien las cumple acumula confianza y respeto. Quien las rompe va perdiendo gente. Hoy vas a aprender 10 acuerdos básicos. Son pocos, son claros y van a servirte en cada chat de Roblox, cada grupo de WhatsApp, cada comentario de TikTok, cada correo institucional, durante toda tu vida en línea.
\end{softbox}

\begin{softbox}{milcMostaza}{milcGris}{Pregunta-puente}
\textit{Cuando alguien escribe \textbf{``ESTÁS TODO LOCO''} en mayúsculas en un chat de juego, ¿\textbf{cómo te sientes}? ¿Sientes que te están hablando bonito o sientes que te están gritando?}
\end{softbox}

\begin{softbox}{milcVerde}{milcGris}{Saber y saber hacer}
Al terminar la clase: (1) podrás \textbf{identificar} los 10 acuerdos de netiqueta; (2) sabrás \textbf{explicar} con tus palabras por qué cada uno cuida la convivencia; (3) habrás \textbf{evaluado} 6 mensajes reales y dirás si cumplen o rompen los acuerdos; (4) habrás \textbf{creado} tu cartel personal con tus propios ejemplos.
\end{softbox}



\small
\begin{tabularx}{\textwidth}{>{\bfseries}p{3.0cm}Y}
\rowcolor{milcGradePrimary!12}Autor & Dr. Álvaro Cárdenas Orozco\\
DBA & Reconoce principios y conceptos propios de la tecnología, relacionando artefactos con su utilización segura (MEN, T\&I 6°)\\
Referentes & Saber ancestral del pregonero · Pensamiento computacional inicial · Cuidado de la palabra\\
Producto final & Un \textbf{cartel personal de los 10 acuerdos de netiqueta} dibujado en hoja completa del cuaderno. Cada acuerdo va con \textbf{un ejemplo de cumplirlo} y \textbf{un ejemplo de romperlo} con tus propias palabras. Más \textbf{una promesa firmada} al final: \emph{``Me comprometo a respetar estos 10 acuerdos en todos los chats, juegos y redes en que esté.''} Lo muestras antes de salir.\\
\end{tabularx}
\normalsize

\puente{Hoy haces 4 cosas: lees situaciones reales, descubres los 10 acuerdos, evalúas mensajes, y haces tu cartel.}

\begin{titlebox}{milcGradeDark}
Ruta MILC: así vamos a aprender
\end{titlebox}
\begin{tabularx}{\textwidth}{>{\centering\arraybackslash}X>{\centering\arraybackslash}X>{\centering\arraybackslash}X>{\centering\arraybackslash}X}
\phasecard{1}{milcVerde}{milcGris}{Escucha\\[-1mm]\small Observo, escucho y pregunto.} &
\phasecard{2}{milcTurquesa}{milcGris}{Sistematización\\[-1mm]\small Pensamiento computacional.} &
\phasecard{3}{milcMagenta}{milcGris}{Praxis\\[-1mm]\small Construyo y aplico.} &
\phasecard{4}{milcVino}{milcGris}{Evaluación\\[-1mm]\small Reflexión liberadora.}
\end{tabularx}


\puente{Empezamos viendo 3 situaciones de chats que tal vez te suenan familiares.}

\begin{titlebox}{milcVerde}
Fase 1 · Escucha
\end{titlebox}

\begin{softbox}{milcVerde}{milcGris}{Escena para escuchar}
\textbf{Actividad 1 · ANALIZA --- Tres situaciones que ya viste} (10 min · individual). Lee estas 3 escenas y en el cuaderno escribe \textbf{qué se hizo mal} en cada una. No tienes que saber los nombres técnicos; cuéntalo con tus palabras. \textbf{Escena A:} En un grupo de WhatsApp del salón, alguien escribe a las 11 de la noche: \emph{``oigan quien tiene la tarea de mate'' --- ``oigan respondan'' --- ``oigan estoy preguntando algo''} (tres mensajes seguidos). \textbf{Escena B:} En Roblox, un jugador escribe en el chat: \emph{``ERES UN MAL JUGADOR APRENDE A JUGAR INÚTIL''}. \textbf{Escena C:} En un grupo de TikTok, alguien comparte una foto de un compañero del salón que se cayó, sin pedirle permiso, con risas.
\end{softbox}

\checkbox Leí las 3 escenas con calma.\\
\checkbox Identifiqué qué se hizo mal en cada una.\\
\checkbox Lo escribí con mis palabras, no copié de otro lado.

\begin{infoband}{milcVerde}


\textbf{Lo que va al cuaderno (Actividad 1 · ANALIZA).} Título: «Actividad 1 --- Tres escenas para mirar». \textbf{Formato:} 3 párrafos cortos, uno por escena. Cada párrafo dice qué se hizo mal y qué hubieras hecho tú. \textbf{Extensión:} 3 párrafos de 3-4 líneas cada uno. \textbf{Sabes que terminaste cuando} si un compañero leyera tu cuaderno, entendería tu opinión sin que tú la explicaras.
\end{infoband}

\puente{Después de las situaciones, aprendes los 10 acuerdos exactos con su nombre.}

\begin{titlebox}{milcTurquesa}
Fase 2 · Sistematización con pensamiento computacional
\end{titlebox}

\begin{softbox}{milcTurquesa}{milcGris}{Cuatro pilares aplicados al tema}
Las 3 escenas que acabaste de leer rompen 3 acuerdos distintos. Pero hay más. En total son \textbf{10 acuerdos básicos de netiqueta} que todo estudiante de tu edad debería conocer. Los vas a aprender ahora.
\end{softbox}

\small
\begin{tabularx}{\textwidth}{>{\bfseries\RaggedRight\arraybackslash}p{3.6cm}Y}
\rowcolor{milcTurquesa!90}\color{white}Pilar & \color{white}\bfseries Aplicación al tema\\
1. Descomposición & \textbf{1. Saluda y despídete.} Como en la vida real: empiezas con un saludo (\emph{``Buenos días''}, \emph{``Hola''}) y terminas con una despedida (\emph{``Gracias''}, \emph{``Hasta mañana''}). No mandes mensajes pelados. \textbf{2. No grites con mayúsculas.} ESCRIBIR TODO EN MAYÚSCULAS es gritar. Solo se usa para una palabra que necesita énfasis, no para frases enteras. \textbf{3. No spam.} No mandes el mismo mensaje 5 veces seguidas. Si nadie responde, espera; si es urgente, llama por teléfono.\\
2. Reconocimiento de patrones & \textbf{4. Cuida los horarios.} No mandes mensajes a horas raras (11 de la noche a 7 de la mañana) salvo emergencia real. Hay personas durmiendo. \textbf{5. Respeta lo que se dice en privado.} Si alguien te cuenta algo en chat privado, no lo compartas con el grupo sin permiso. \textbf{6. No publiques fotos de otros sin pedirles permiso.} Aunque la foto sea graciosa. Aunque la persona ``no se vaya a dar cuenta''. Pedir permiso es regla básica.\\
3. Abstracción & \textbf{7. Piensa antes de escribir.} En la vida real, las palabras se las lleva el viento. En internet, las palabras quedan. Antes de enviar algo molesto o brusco, lee otra vez: \emph{``¿De verdad quiero decir esto así?''}. \textbf{8. No insultes ni humilles, ni en broma.} Lo que en persona puede ser broma, en pantalla suele leerse como insulto. No hay tono de voz, no hay sonrisa. Mejor evítalo.\\
4. Algoritmo conceptual & \textbf{9. Si algo está mal, no lo reenvíes.} Cadenas falsas, fotos íntimas de alguien, noticias raras, retos peligrosos. \textbf{Reenviarlo te hace parte del problema}. La regla es: si lo dudas, no lo mandes. \textbf{10. Habla con un adulto cuando te asustes.} Si alguien en internet te molesta, te amenaza, te pide fotos raras, te dice cosas que te hacen sentir mal --- cuéntale a un adulto de confianza (mamá, papá, profe, abuela). No estás solo.\\
\end{tabularx}
\normalsize

\begin{drawbox}{milcTurquesa}{Los 10 acuerdos de netiqueta}
\textbf{1.} Saluda y despídete. \\[1mm]
\textbf{2.} No grites con mayúsculas. \\[1mm]
\textbf{3.} No spam (no mandes lo mismo 5 veces). \\[1mm]
\textbf{4.} Cuida los horarios (nada de mensajes a las 11pm). \\[1mm]
\textbf{5.} Lo privado se queda privado. \\[1mm]
\textbf{6.} No publiques fotos de otros sin permiso. \\[1mm]
\textbf{7.} Piensa antes de escribir. \\[1mm]
\textbf{8.} No insultes ni humilles, ni en broma. \\[1mm]
\textbf{9.} Si algo está mal, no lo reenvíes. \\[1mm]
\textbf{10.} Cuéntale a un adulto cuando algo te asuste.
\end{drawbox}

\begin{softbox}{milcMostaza}{milcGris}{Errores comunes y su corrección}
\textbf{Mitos sobre netiqueta que vas a oír de tus compañeros.} (1) \emph{``Si no contesta es porque es altanero''} --- No. Hay clases, vidas, sueño. Esperar es respeto. (2) \emph{``Si no escribo en mayúsculas no se nota lo que digo''} --- Sí se nota. Usa una palabra puntual en mayúscula, no frases enteras. (3) \emph{``Si yo lo recibí, lo puedo reenviar''} --- Falso. Solo porque te lo mandaron no significa que ya es público. (4) \emph{``Si alguien me molesta, le respondo con otro insulto''} --- Eso enciende más el problema. La salida es bloquear o avisar a un adulto. (5) \emph{``Es solo internet, no es real''} --- Es muy real. Lo que pasa en pantalla afecta cómo te sientes y cómo te ven.
\end{softbox}

\begin{infoband}{milcTurquesa}


\textbf{Lo que va al cuaderno (Actividad 2 · EXPLICA).} Título: «Actividad 2 --- Los 10 acuerdos con mis palabras». \textbf{Formato:} lista numerada del 1 al 10, cada acuerdo en una sola línea, escrito con \emph{tu propia voz}, no copiado. \textbf{Extensión:} 10 líneas. \textbf{Sabes que terminaste cuando} podrías recitar los 10 sin abrir el cuaderno.
\end{infoband}

\puente{Con los 10 acuerdos en mente, dibujas tu cartel personal en el cuaderno.}

\begin{titlebox}{milcMagenta}
Fase 3 · Praxis con pensamiento computacional
\end{titlebox}

\begin{softbox}{milcMagenta}{milcGris}{Construyendo el producto con los cuatro pilares}
\textbf{Actividad 3 · CREA --- Tu cartel personal de netiqueta} (25 min · individual). Vas a hacer en una hoja completa del cuaderno un \textbf{cartel personal} con los 10 acuerdos. Cada acuerdo lleva 2 ejemplos pequeños: uno de \textbf{cumplir} y uno de \textbf{romper} el acuerdo. Vas a inventar los ejemplos tú; no los copies del libro. Al final, firmas el cartel con tu nombre como una promesa.
\end{softbox}

\small
\begin{tabularx}{\textwidth}{>{\bfseries\RaggedRight\arraybackslash}p{3.6cm}Y}
\rowcolor{milcMagenta!90}\color{white}Pilar & \color{white}\bfseries Aplicación al producto\\
Descomposición & \textbf{Paso 1.} En la parte superior de la hoja escribe el título grande: \textbf{``Mis 10 acuerdos de netiqueta''} y debajo: \emph{``Promesa de [tu nombre], grado 6[tu letra], año 2026''}. Subraya el título.\\
Patrones & \textbf{Paso 2.} Divide la hoja en \textbf{10 cuadritos} (2 columnas, 5 filas, por ejemplo) o haz una lista bien ordenada. Cada cuadrito o renglón es para 1 acuerdo. Numéralos del 1 al 10.\\
Abstracción & \textbf{Paso 3.} En cada cuadrito escribe: arriba el \textbf{nombre del acuerdo} (corto, máximo 5 palabras). Abajo \textbf{un ejemplo de cumplirlo} (\emph{``Cumpliendo: ...''}) y \textbf{un ejemplo de romperlo} (\emph{``Rompiendo: ...''}). Los ejemplos los inventas tú, deben ser cosas que pasan en tu vida.\\
Algoritmo de armado & \textbf{Paso 4.} Para los acuerdos que tienen que ver con grupos del colegio (por ejemplo el acuerdo 3 o el 6), trata de que el ejemplo sea del salón de 6° o de una situación que ya viste o que podría pasar. Esto hace tu cartel \emph{tuyo}.\\
\end{tabularx}
\normalsize

\begin{drawbox}{milcMagenta}{Antes de mostrar tu cartel}
\checkbox El cartel ocupa una hoja entera del cuaderno. \\[1mm]
\checkbox Están los 10 acuerdos numerados del 1 al 10. \\[1mm]
\checkbox Cada acuerdo tiene un ejemplo de cumplir y uno de romper. \\[1mm]
\checkbox Los ejemplos los inventaste tú (no son copias). \\[1mm]
\checkbox Cada acuerdo tiene un ícono pequeño. \\[1mm]
\checkbox La promesa está firmada con nombre y fecha.
\end{drawbox}

\begin{softbox}{milcMagenta}{milcGris}{Plantilla de trabajo}
\textbf{Plantilla de cada cuadrito.} \textit{Número}: 1 al 10. \textit{Nombre del acuerdo}: máximo 5 palabras. \textit{Cumpliendo}: 1 frase. \textit{Rompiendo}: 1 frase. \textit{Ícono pequeño}: el que tú quieras.
\end{softbox}

\begin{infoband}{milcMagenta}


\textbf{Lo que va al cuaderno (Actividad 3 · CREA).} Título: «Actividad 3 --- Mis 10 acuerdos de netiqueta». \textbf{Formato:} hoja completa del cuaderno con 10 cuadritos, ejemplos y firma. \textbf{Extensión:} 1 hoja entera. \textbf{Sabes que terminaste cuando} cualquier compañero de otro salón podría leerlo y entenderlo sin que tú lo expliques.
\end{infoband}

\puente{El producto es el cartel firmado: lo guardas como compromiso tuyo.}

\begin{titlebox}{milcMostaza}
Producto de la sesión
\end{titlebox}
\textbf{Cartel personal de los 10 acuerdos de netiqueta · firmado}.
\begin{softbox}{milcMostaza}{milcGris}{Criterios de calidad}
(1) El cartel ocupa una página completa. (2) Están los 10 acuerdos en orden del 1 al 10. (3) Cada acuerdo tiene su ejemplo de cumplirlo y de romperlo. (4) Los ejemplos son inventados por el estudiante, contextualizados. (5) El cartel está firmado y con fecha.
\end{softbox}

\puente{Antes de salir, revisas que tu cartel tiene los 10 acuerdos, los ejemplos y la firma.}

\begin{titlebox}{milcVino}
Fase 4 · Evaluación liberadora — cinco dimensiones
\end{titlebox}

\begin{softbox}{milcVino}{milcGris}{Sentido de la evaluación}
Evaluar no es solo revisar una nota. Es reconocer qué aprendí, qué cambió en mí, qué emociones manejé, cómo me relaciono con la ciudad y la comunidad, y qué le dejaré a quien viene detrás.
\end{softbox}

\textbf{Mi reflexión liberadora en cinco dimensiones.} Completa cada frase iniciada:

\small
\begin{tabularx}{\textwidth}{>{\bfseries\RaggedRight\arraybackslash}p{3.4cm}Y>{\RaggedRight\arraybackslash}p{4.6cm}}
\rowcolor{milcVino}\color{white}Dimensión & \color{white}\bfseries Mi respuesta (frame starter) & \color{white}\bfseries Algo que descubrí\\
1. Personal & Lo que más cambió en mí fue \linead{6.5cm} & \linead{4.2cm}\\
2. Emocional & La emoción más fuerte fue \linead{2.5cm} y la regulé \linead{3cm} & \linead{4.2cm}\\
3. Ciudadana & En democracia esto importa porque \linead{6cm} & \linead{4.2cm}\\
4. Local (Cartago) & En mi barrio sería útil para \linead{6.5cm} & \linead{4.2cm}\\
5. Intergeneracional & A mi abuela/abuelo le diría \linead{6.5cm} & \linead{4.2cm}\\
\end{tabularx}
\normalsize

\begin{infoband}{milcVino}
\textbf{Para la próxima sustentación.} Vuelve a esta tabla en seis meses. Verás que las respuestas cambian — no porque lo aprendido cambie, sino porque tú cambias. La evaluación liberadora es una conversación contigo a través del tiempo.
\end{infoband}


\puente{Cerramos con 3 voces que te ayudan a entender por qué el respeto pesa más en pantalla que cara a cara.}

\begin{titlebox}{milcGradeDark}
Triángulo de pensamiento · cierre filosófico
\end{titlebox}

\begin{softbox}{milcVino}{milcGris}{Dussel · lente del nosotros}
\textit{````Antes de hablar a otro, primero hay que reconocer que el otro es persona como tú.''''}

En internet no ves la cara del otro. No ves si está cansado, si llora, si tiene un mal día. \textbf{Por eso la regla es más fuerte ahí}: tratar al otro como persona, no como avatar. Cada nombre que aparece en un chat es un ser humano del otro lado. Saludarlo, escribirle con cuidado, no humillarlo --- eso es reconocer a la persona.

\textbf{Cuando escribo en línea, ¿\textbf{recuerdo que el del otro lado es una persona} o trato a las personas como si fueran nombres en una pantalla?}
\end{softbox}
\linead{\linewidth}\\[1mm]
\linead{\linewidth}

\begin{softbox}{milcMostaza}{milcGris}{Estoicismo · Epicteto · cuidado interior}
\textit{````No es lo que te dicen lo que te hace daño, es cómo respondes a lo que te dicen.''''}

Si alguien te escribe algo molesto en un chat, lo que te enciende es \emph{cómo reaccionas tú}, no lo que dijo el otro. \textbf{Tienes 3 caminos:} (1) responder con otro insulto y encender el problema, (2) ignorar y bloquear, (3) avisarle a un adulto si es serio. Los caminos 2 y 3 son los inteligentes. La rabia se enfría con un vaso de agua, no escribiendo.

\textbf{Cuando algo me molesta en línea, ¿\textbf{respondo de una vez} o tomo un momento antes de escribir?}
\end{softbox}
\linead{\linewidth}\\[1mm]
\linead{\linewidth}

\begin{softbox}{milcTurquesa}{milcGris}{Floridi · lente de la infoesfera}
\textit{````Cada mensaje que envías deja huella. La gente recordará no lo que dijiste, sino cómo lo dijiste.''''}

De todos los mensajes que envíes en tu vida, \textbf{la gente se acordará más del tono que del contenido}. Pueden no recordar exactamente qué pediste, pero recordarán si pediste con respeto o si llegaste empujando. La netiqueta no es debilidad: es saber que el tono pesa.

\textbf{Si la gente solo recordara cómo me hablo en internet, ¿\textbf{qué pensarían} de mí?}
\end{softbox}
\linead{\linewidth}\\[1mm]
\linead{\linewidth}

\begin{titlebox}{milcVino}
Compromiso final
\end{titlebox}

\small
\begin{tabularx}{\textwidth}{>{\RaggedRight\arraybackslash}p{3.0cm}YYY}
\rowcolor{milcVino}\color{white}\bfseries Criterio & \color{white}\bfseries Lo logré & \color{white}\bfseries En proceso & \color{white}\bfseries Necesito apoyo\\
Comprensión del tema & Explico ideas centrales con mis palabras. & Reconozco ideas pero falta precisión. & Necesito repasar conceptos base.\\
Pensamiento computacional & Apliqué los 4 pilares al diseñar mi producto. & Apliqué algunos pilares. & Necesito releer la fase 2.\\
Aplicación y producto & Mi producto cumple los criterios y se puede revisar. & Producto funcional con ajustes pendientes. & Aún en construcción.\\
Reflexión 5 dimensiones & Respondí cada dimensión con honestidad. & Respondí algunas con detalle. & Mis respuestas son muy generales.\\
Triángulo de pensamiento & Conecté las 3 voces con mi trabajo real. & Conecté al menos una. & No relaciono las voces con mi proceso.\\
\end{tabularx}
\normalsize

\textbf{Hoy entendí que:}\\[1mm]
\linead{\linewidth}\\[1mm]
\linead{\linewidth}

\textbf{Mi compromiso para la próxima sesión es:}\\[1mm]
\linead{\linewidth}\\[1mm]
\linead{\linewidth}

\begin{softbox}{milcMostaza}{milcGris}{Cita de despedida · Marco Aurelio}
\textit{``El obstáculo en el camino se vuelve el camino. Lo que estorba al camino se vuelve camino.''}\\[1mm]
Cada error de hoy, bien observado, se convierte en pieza del camino que pisarás mañana.
\end{softbox}

\small
\begin{tabularx}{\textwidth}{>{\bfseries}p{3.6cm}Y}
\rowcolor{milcGradeDark!12}Nombre completo: & \linead{10cm}\\
Fecha: & \linead{4cm} \quad Curso/grupo: \linead{4cm}\\
Mi firma: & \linead{9cm}\\
\end{tabularx}
\normalsize

\begin{infoband}{milcGradeDark}
\textbf{Para la próxima sesión.} Guarda tu cartel en el cuaderno. Esta semana, cada vez que entres a un chat, mira mentalmente uno de los 10 acuerdos y pregúntate: \emph{``¿Lo estoy cumpliendo?''}. La próxima clase vamos a viajar al pasado: vas a descubrir \textbf{cómo eran los medios de comunicación} antes de internet --- desde el pregonero hasta la radio y la televisión --- y por qué esta historia te ayuda a entender mejor TikTok y WhatsApp.
\end{infoband}

\end{document}
